\section{POSIX Shell Automation: Administrative Scripts}

The Unix philosophy dictates that tools should "do one thing and do it well." The five scripts
developed for this audit demonstrate practical Linux administration and shell scripting concepts,
from reading system identity to interactive user-driven file generation. Each script is
documented below with its full source code, an explanation of its purpose, and the shell
constructs it employs.

% ------------------------------------------------------------------
\subsection{Script 1 — System Identity Report}
% ------------------------------------------------------------------

\textbf{Purpose:} Acts as an administrative welcome screen for the audit environment. The script
dynamically collects kernel, user, uptime, and distribution information, then displays it in a
structured format. It concludes with a message identifying the GPLv2 license covering the
underlying Linux OS — directly relevant to the open-source audit context.

\textbf{Shell Concepts Used:} Variables, command substitution \texttt{\$(...)}, \texttt{echo}
with formatted output, and reading structured files with \texttt{grep} and \texttt{cut}.

\begin{lstlisting}[language=bash, caption={01\_system\_identity.sh — System Identity Report}]
#!/bin/bash
# Script 1: System Identity Report
# Author: Ekjot Singh | Roll: 24BCE10284 | Course: Open Source Software
# Description: Displays a welcome screen with system identity
#              and OS license information.
# Concepts: Variables, command substitution $(), echo, output formatting.

# --- Variables: storing static audit metadata ---
STUDENT_NAME="Ekjot Singh"
SOFTWARE_CHOICE="Git"

# --- Command substitution: dynamically fetching live system data ---
KERNEL=$(uname -r)
USER_NAME=$(whoami)
USER_HOME=$HOME
UPTIME=$(uptime -p)
CURRENT_DATE=$(date +"%A, %B %d, %Y %T")

# Reading the distribution name from the standard os-release file.
# grep finds the PRETTY_NAME line; cut extracts the value between quotes.
DISTRO=$(grep '^PRETTY_NAME' /etc/os-release | cut -d'"' -f2)

# --- Display: formatted output using echo ---
echo "=================================================="
echo " Open Source Audit - $STUDENT_NAME"
echo " Focus Software: $SOFTWARE_CHOICE"
echo "=================================================="
echo "Date/Time  : $CURRENT_DATE"
echo "Distro     : $DISTRO"
echo "Kernel     : $KERNEL"
echo "User       : $USER_NAME"
echo "Home Dir   : $USER_HOME"
echo "Uptime     : $UPTIME"
echo "--------------------------------------------------"
# Explicitly stating the open-source license governing the OS
echo "License    : The Linux kernel is released under the"
echo "             GNU General Public License v2 (GPLv2)."
echo "             Source code must remain freely available."
echo "=================================================="
\end{lstlisting}

% ------------------------------------------------------------------
\subsection{Script 2 — FOSS Package Inspector}
% ------------------------------------------------------------------

\textbf{Purpose:} Checks whether the chosen software (Git) is installed on the system,
retrieves its version and description using the \texttt{dpkg} package manager, and uses a
\texttt{case} statement to print a philosophical note about the package's role in the
open-source ecosystem.

\textbf{Shell Concepts Used:} \texttt{if-then-else} for conditional logic, \texttt{case}
statement for multi-branch decisions, \texttt{dpkg} with \texttt{grep -E} for structured
package queries, and output redirection with \texttt{\&>/dev/null} to suppress errors.

\begin{lstlisting}[language=bash, caption={02\_package\_inspector.sh — FOSS Package Inspector}]
#!/bin/bash
# Script 2: FOSS Package Inspector
# Author: Ekjot Singh | Roll: 24BCE10284
# Description: Inspects whether the target FOSS package is installed,
#              retrieves its version, and prints its ecosystem philosophy.
# Concepts: if-then-else, case statement, dpkg, grep -E pipe.

PACKAGE="git"

echo "Inspecting system for FOSS package: $PACKAGE..."
echo "------------------------------------------------"

# if-then-else: check installation status.
# dpkg -s returns exit code 0 if installed, non-zero if not.
# &>/dev/null redirects both stdout and stderr to suppress noisy output.
if dpkg -s "$PACKAGE" &>/dev/null; then
    echo "[SUCCESS] $PACKAGE is installed on this system."
    # Pipe dpkg output through grep with extended regex to extract
    # Version and Description fields only.
    dpkg -s "$PACKAGE" | grep -E '^Version|^Description'
else
    echo "[WARNING] $PACKAGE is NOT installed."
    echo "          Install it with: sudo apt install $PACKAGE"
fi

echo "------------------------------------------------"

# case statement: print a philosophy note based on the package name.
# This demonstrates multi-branch decision making without chained if-elif.
case $PACKAGE in
    git)
        echo "Git: The decentralized engine of global collaboration" \
             "and code sovereignty."
        ;;
    apache2|httpd)
        echo "Apache: The web server that built the open internet."
        ;;
    mysql|mariadb-server)
        echo "MySQL: Open source at the heart of millions of apps."
        ;;
    vlc)
        echo "VLC: The universal media player — free from proprietary codecs."
        ;;
    *)
        echo "An open-source tool contributing to the vast Linux ecosystem."
        ;;
esac
\end{lstlisting}

% ------------------------------------------------------------------
\subsection{Script 3 — Disk and Permission Auditor}
% ------------------------------------------------------------------

\textbf{Purpose:} Iterates through a predefined list of critical Linux directories and
reports their POSIX permissions, ownership, and disk usage. A separate block then checks
for Git's specific system-level configuration file. This script is directly relevant to the
Linux Footprint section of the audit (Part B).

\textbf{Shell Concepts Used:} \texttt{for} loop over an array, \texttt{[ -d ]} and
\texttt{[ -f ]} conditional tests, \texttt{du -sh} for disk usage, \texttt{ls -ld} with
\texttt{awk} to extract permission fields.

\begin{lstlisting}[language=bash, caption={03\_disk\_auditor.sh — Disk and Permission Auditor}]
#!/bin/bash
# Script 3: Disk and Permission Auditor
# Author: Ekjot Singh | Roll: 24BCE10284
# Description: Audits key Linux directories for permissions, ownership,
#              and disk usage. Also checks Git's config footprint.
# Concepts: for loop, arrays, [ -d ] test, du, ls -ld, awk, cut.

# Define the audited software for labelling purposes
SOFTWARE_CHOICE="Git"

# Array of critical system directories to audit
DIRS=("/etc" "/var/log" "/home" "/usr/bin" "/tmp")

echo "========================================"
echo " Directory Audit Report"
echo "========================================"

# for loop: iterate through every directory in the array
for DIR in "${DIRS[@]}"; do
    # Conditional test: [ -d ] returns true if path is a directory
    if [ -d "$DIR" ]; then
        # awk extracts columns 1 (permissions), 3 (owner), 4 (group)
        # from ls -ld long-format output
        PERMS=$(ls -ld "$DIR" | awk '{print $1, $3, $4}')
        # du -sh gives human-readable size; cut -f1 takes just the number
        SIZE=$(du -sh "$DIR" 2>/dev/null | cut -f1)
        echo "$DIR"
        echo "  Permissions/Owner : $PERMS"
        echo "  Disk Usage        : $SIZE"
        echo "----------------------------------------"
    else
        echo "$DIR — does not exist on this system"
        echo "----------------------------------------"
    fi
done

echo ""
echo "Checking $SOFTWARE_CHOICE system configuration footprint..."
echo "========================================"

# [ -f ] tests for a regular file (not a directory)
GIT_CONFIG="/etc/gitconfig"
GIT_BIN="/usr/bin/git"

if [ -f "$GIT_CONFIG" ]; then
    GIT_PERMS=$(ls -ld "$GIT_CONFIG" | awk '{print $1, $3, $4}')
    echo "$GIT_CONFIG => $GIT_PERMS"
else
    echo "Global Git config ($GIT_CONFIG) not found on this system."
fi

if [ -f "$GIT_BIN" ]; then
    BIN_PERMS=$(ls -ld "$GIT_BIN" | awk '{print $1, $3, $4}')
    echo "$GIT_BIN => $BIN_PERMS"
else
    echo "Git binary ($GIT_BIN) not found."
fi
\end{lstlisting}

% ------------------------------------------------------------------
\subsection{Script 4 — Log File Analyzer}
% ------------------------------------------------------------------

\textbf{Purpose:} Reads a system log file line by line, counts how many lines contain a
specified keyword (defaulting to \texttt{error}), and prints a summary alongside the five
most recent matching lines. It includes a do-while style retry loop that prompts the user
for a valid file path if the original argument is missing or empty.

\textbf{Shell Concepts Used:} \texttt{while read} loop for safe line-by-line file parsing,
\texttt{if-then} inside a loop, counter variables with arithmetic expansion
\texttt{\$((COUNT + 1))}, command-line positional arguments \texttt{\$1}/\texttt{\$2} with
default values, and \texttt{tail | grep} pipeline for extracting recent anomalies.

\begin{lstlisting}[language=bash, caption={04\_log\_analyzer.sh — Log File Analyzer}]
#!/bin/bash
# Script 4: Log File Analyzer
# Author: Ekjot Singh | Roll: 24BCE10284
# Usage: ./04_log_analyzer.sh /var/log/syslog error
# Description: Scans a log file for a keyword, counts occurrences,
#              and extracts the 5 most recent matching lines.
# Concepts: while-read loop, if-then, counter vars, $1/$2 args,
#           do-while style retry, tail+grep pipeline.

# Positional arguments with default fallback values using :-
LOGFILE=${1:-"/var/log/syslog"}
KEYWORD=${2:-"error"}   # Default keyword if none provided
COUNT=0                 # Counter variable, initialised to zero

# Do-while style retry: keep prompting until a valid non-empty file is given.
# [ ! -s ] is true when the file does not exist OR is empty (size = 0).
while [ ! -s "$LOGFILE" ]; do
    echo "Error: File '$LOGFILE' is missing or empty."
    read -p "Please enter a valid log file path: " LOGFILE
done

echo "Scanning '$LOGFILE' for keyword: '$KEYWORD'..."
echo "------------------------------------------------"

# while IFS= read -r LINE: reads the file line by line safely,
# preserving whitespace (IFS=) and backslashes (-r).
while IFS= read -r LINE; do
    # if-then: check each line for keyword match (case-insensitive with -i)
    if echo "$LINE" | grep -iq "$KEYWORD"; then
        # Arithmetic expansion increments the counter
        COUNT=$((COUNT + 1))
    fi
done < "$LOGFILE"

echo "Keyword '$KEYWORD' found $COUNT times in $LOGFILE"
echo "------------------------------------------------"

# Print last 5 matching lines using tail + grep pipeline
if [ "$COUNT" -gt 0 ]; then
    echo "Last 5 occurrences:"
    grep -i "$KEYWORD" "$LOGFILE" | tail -n 5
else
    echo "No occurrences of '$KEYWORD' found."
fi
\end{lstlisting}

% ------------------------------------------------------------------
\subsection{Script 5 — Open Source Manifesto Generator}
% ------------------------------------------------------------------

\textbf{Purpose:} An interactive script that collects three user responses about open-source
philosophy and composes a personalised manifesto paragraph, saving it to a \texttt{.txt}
file. It demonstrates user input, string concatenation, file redirection, and the
\texttt{alias} concept relevant to Unit 5 of the course.

\textbf{Shell Concepts Used:} \texttt{read} for interactive user input, string concatenation
via variable expansion, file creation/overwrite with \texttt{>} and append with \texttt{>>},
\texttt{date} command substitution, and a live \texttt{alias} declaration demonstrating
the alias concept.

\begin{lstlisting}[language=bash, caption={05\_manifesto\_generator.sh — Open Source Manifesto Generator}]
#!/bin/bash
# Script 5: Open Source Manifesto Generator
# Author: Ekjot Singh | Roll: 24BCE10284
# Description: Interactively generates a personalised open-source
#              philosophy manifesto and saves it to a .txt file.
# Concepts: read (user input), string concatenation, > and >> redirection,
#           date command, alias declaration.

# --- Alias concept demonstration ---
# An alias maps a long command to a short custom label.
# Here we define 'today' as a short alias for a formatted date command.
# In a login shell, aliases are typically stored in ~/.bashrc or ~/.bash_aliases.
alias today='date "+%d %B %Y"'

echo "=================================================="
echo " Open Source Manifesto Generator"
echo "=================================================="
echo "Answer three questions to generate your manifesto."
echo ""

# read -p: displays the prompt and waits for user input.
# The response is stored in the named variable.
read -p "1. Name one open-source tool you use every day: " TOOL
read -p "2. In one word, what does 'freedom' mean to you? " FREEDOM
read -p "3. Name one thing you would build and share freely: " BUILD

# Command substitution: capture the current date into a variable
DATE=$(date '+%d %B %Y')

# Output filename includes the current username for personalisation
OUTPUT="manifesto_$(whoami).txt"

echo ""
echo "Composing your manifesto..."
echo "--------------------------------------------------"

# > operator: creates the file (or overwrites it if it already exists)
# Writing the header block line by line
echo "=== My Open Source Manifesto ===" > "$OUTPUT"
echo "Author : $(whoami)"               >> "$OUTPUT"
echo "Date   : $DATE"                   >> "$OUTPUT"
echo ""                                 >> "$OUTPUT"

# >> operator: appends the manifesto paragraph to the file without overwriting.
# String concatenation is achieved by embedding variables directly
# inside a double-quoted echo string.
echo "I believe that software should empower rather than restrict. \
The tool I rely on daily, $TOOL, represents this philosophy perfectly. \
To me, open source is synonymous with $FREEDOM — it ensures that code \
remains a public good, not a locked box. Standing on the shoulders of \
the community, I pledge to build $BUILD and share it freely, so others \
may learn from my work just as I have learned from theirs." >> "$OUTPUT"

echo "--------------------------------------------------"
echo "Manifesto saved to: $OUTPUT"
echo "--------------------------------------------------"

# Display the completed file using cat
cat "$OUTPUT"
\end{lstlisting}
